\chapter{Virial identities}
\label{chap:virial}

A \emph{virial identity} pairs the eigenvalue equation with a well chosen
first order operator $D$ and uses the commutator $[H,D]$. All three sectors
use the same mechanism with different vector fields. We first record the
algebra for smooth functions, where the identity takes three lines, and then
the version on the original $L^2$ domain, which is where the analytic work
lies.

\begin{definition}[Adapted vector field]
\label{def:adapted-field}
\uses{def:dipole, def:lift-operators}
\lean{NoEmbeddedEigenvalues.IsAdaptedOn}
\leanok
Let $Y=\nabla^\perp\Psi=(\sin y,-\sin x)$, so that $L_0=Y\cdot\nabla$. A real
vector field $W=(W_1,W_2)$ on a subset $U$ of $\T^2$ (or of its lift $\R^2$) is
\emph{adapted} with \emph{commutator coefficient} $f$ if the Lie bracket
$[Y,W]=(Y\cdot\nabla)W-(W\cdot\nabla)Y$ equals $fY$ on $U$. In components,
\[
  L_0W_1-\cos y\,W_2=f\sin y,\qquad L_0W_2+\cos x\,W_1=-f\sin x .
\]
For $W\in C^1$ this is the operator identity $[L_0,W\cdot\nabla]=f\,L_0$ on
$C^2$ functions (Lemma~\ref{lem:commutator-pointwise}).
\end{definition}

\begin{definition}[Half-density generator]
\label{def:half-density-generator}
\uses{def:adapted-field, def:lift-operators}
\lean{NoEmbeddedEigenvalues.fieldDeriv,
  NoEmbeddedEigenvalues.divergence,
  NoEmbeddedEigenvalues.halfDensityGen}
\leanok
For a $C^1$ vector field $W$ on $\T^2$ (or on its lift), its
\emph{half-density generator} is $D=W\cdot\nabla+\tfrac12\Div W$, where
$W\cdot\nabla u=W_1\partial_xu+W_2\partial_yu$ and
$\Div W=\partial_xW_1+\partial_yW_2$.
\end{definition}

\begin{lemma}[Pointwise commutator identities]
\label{lem:commutator-pointwise}
\uses{def:adapted-field, def:half-density-generator}
\lean{NoEmbeddedEigenvalues.IsAdaptedOn.classicalL0_fieldDeriv,
  NoEmbeddedEigenvalues.IsAdaptedOn.classicalL0_divergence,
  NoEmbeddedEigenvalues.IsAdaptedOn.classicalL0_halfDensityGen}
\leanok
Let $U$ be open, $W$ of class $C^2$ and adapted on $U$ with a differentiable
coefficient $f$, and $d=\Div W$. Then on $U$, for every $C^2$ function $u$,
\begin{enumerate}
  \item $[L_0,W\cdot\nabla]u=f\,L_0u$;
  \item $L_0d=L_0f$;
  \item $[L_0,D]u=\tfrac12\bigl(fL_0u+L_0(fu)\bigr)$.
\end{enumerate}
Item (1) holds at each point of $U$ and needs only $W$ differentiable.
\end{lemma}

\begin{proof}
\uses{lem:lift-calculus}
\leanok
For (1), $[Y\cdot\nabla,W\cdot\nabla]u=[Y,W]\cdot\nabla u$ because the second
derivatives of $u$ cancel by symmetry; then use $[Y,W]=fY$. For (2), the two
components of $[Y,W]-fY$ vanish on $U$, so their derivatives do; the
$x$-derivative of the first plus the $y$-derivative of the second is
$L_0d-L_0f$, by symmetry of the mixed partials of $W$. (Equivalently: since
$\Div Y=0$, $\Div[Y,W]=Y\cdot\nabla\Div W$, while $\Div(fY)=Y\cdot\nabla f$.)
For (3), $[L_0,D]u=[L_0,W\cdot\nabla]u+\tfrac12(L_0d)u=fL_0u+\tfrac12(L_0f)u
=\tfrac12(fL_0u+L_0(fu))$ by (1) and (2).
\end{proof}

\begin{lemma}[Skew-symmetry]
\label{lem:commutator-algebra}
\uses{def:half-density-generator, def:square-integral}
\lean{NoEmbeddedEigenvalues.squarePairing_halfDensityGen,
  NoEmbeddedEigenvalues.halfDensityGen_mul_conj_add,
  NoEmbeddedEigenvalues.squarePairing_classicalL0,
  NoEmbeddedEigenvalues.halfDensityGen_velocity}
\leanok
Let $W$ be a periodic $C^1$ vector field. For periodic $C^1$ functions $u,v$,
\begin{enumerate}
  \item $D$ is skew-symmetric: $\ip{Du}{v}=-\ip{u}{Dv}$;
  \item $L_0$ is skew-symmetric, so $H=-iL_0$ is symmetric:
        $\ip{L_0u}v=-\ip u{L_0v}$.
\end{enumerate}
\end{lemma}

\begin{proof}
\uses{lem:square-integral-partial, lem:square-pairing, lem:lift-calculus}
\leanok
Pointwise, $(Du)\bar v+u\overline{Dv}=\Div(u\bar vW)$: the half divergence in
$D$ supplies exactly the term $(\Div W)u\bar v$. The integral of a divergence
of a periodic $C^1$ field vanishes. Item (2) is the case $W=Y$, for which
$\Div Y=0$ and $D=L_0$.
\end{proof}

\begin{remark}
Multiplying Lemma~\ref{lem:commutator-pointwise}(3) by $-i$ gives
$[H,D]=\tfrac12(fH+Hf)$, the form used in the manuscript.
\end{remark}

\begin{lemma}[Virial identity, smooth case]
\label{lem:virial-smooth}
\uses{def:adapted-field, def:half-density-generator, def:eigen-problem, def:square-integral}
\lean{NoEmbeddedEigenvalues.virial_smooth}
\leanok
Let $W$ be a periodic $C^2$ vector field, adapted on $\R^2$ with a periodic $C^1$
coefficient $f$. Let $\omega,\psi$ be periodic $C^2$ functions with
$L_0(\omega+\psi)=i\lambda\omega$ and $\lambda\ne0$ (for instance a smooth
eigenfunction and $\psi=\Delta^{-1}\omega$). Then
\begin{equation*}
  \operatorname{Re}\ip{f\omega}{\omega+\psi}=2\operatorname{Re}\ip{D\psi}{\omega}
\end{equation*}
(manuscript \texttt{(ee:commutator-limit)}). The relation $\psi=\Delta^{-1}\omega$
is not needed for this form.
\end{lemma}

\begin{proof}
\uses{lem:commutator-pointwise, lem:commutator-algebra, lem:square-pairing}
\leanok
Put $a=\omega+\psi$, so $Ha=\lambda\omega$. On the one hand
$\ip{[H,D]a}{a}=\tfrac12(\ip{fHa}{a}+\ip{fa}{Ha})
=\lambda\operatorname{Re}\ip{f\omega}{a}$. On the other hand, using
self-adjointness of $H$ and skew-symmetry of $D$,
$\ip{[H,D]a}{a}=\ip{Da}{Ha}-\ip{D\,Ha}{a}
=\lambda\ip{Da}{\omega}+\lambda\ip{\omega}{Da}
=2\lambda\operatorname{Re}\ip{Da}{\omega}$, and
$\operatorname{Re}\ip{D\omega}{\omega}=0$. Divide by $\lambda$.
\end{proof}

\begin{lemma}[Virial identity, smooth case, quadratic form]
\label{lem:virial-smooth-quadratic}
\uses{lem:virial-smooth}
\lean{NoEmbeddedEigenvalues.virial_smooth_quadratic,
  NoEmbeddedEigenvalues.squarePairing_partialX,
  NoEmbeddedEigenvalues.squarePairing_partialY,
  NoEmbeddedEigenvalues.re_squarePairing_smul_partialX,
  NoEmbeddedEigenvalues.re_squarePairing_smul_partialY,
  NoEmbeddedEigenvalues.re_squarePairing_smul_laplacian,
  NoEmbeddedEigenvalues.partialX_halfDensityGen,
  NoEmbeddedEigenvalues.partialY_halfDensityGen,
  NoEmbeddedEigenvalues.re_squarePairing_halfDensityGen_self}
\leanok
In the setting of Lemma~\ref{lem:virial-smooth}, assume moreover $\Delta\psi=\omega$,
with $\psi\in C^4$, $W\in C^3$ and $f\in C^2$. Then
\[
  0=\int_{\T^2}f|\Delta\psi|^2
   +\int_{\T^2}(2\sym DW-fI)\nabla\psi\cdot\overline{\nabla\psi}
   +\frac12\int_{\T^2}\Delta(f-\Div W)|\psi|^2
\]
(manuscript \texttt{(ee:bounded-form)}).
\end{lemma}

\begin{proof}
\uses{lem:commutator-algebra, lem:square-pairing, lem:lift-calculus}
\leanok
Substitute $\omega=\Delta\psi$ in Lemma~\ref{lem:virial-smooth} and use the
integrations by parts
$\operatorname{Re}\int f\Delta\psi\,\bar\psi=-\int f|\nabla\psi|^2
+\tfrac12\int\Delta f|\psi|^2$ and
$2\operatorname{Re}\int D\psi\,\overline{\Delta\psi}
=-2\int\sym DW\nabla\psi\cdot\overline{\nabla\psi}+\tfrac12\int\Delta(\Div W)|\psi|^2$.

All integrations by parts follow from skew-symmetry
(Lemma~\ref{lem:commutator-algebra}): $\partial_x,\partial_y$ are the generators $D$
of the constant fields $(1,0)$, $(0,1)$. For real $g$ this gives
$\operatorname{Re}\ip{gu}{\partial_ju}=-\tfrac12\ip{(\partial_jg)u}{u}$. For the second
identity, $\partial_j(D\psi)=D(\partial_j\psi)+(\partial_jW)\cdot\nabla\psi
+\tfrac12(\partial_j\Div W)\psi$ and $\operatorname{Re}\ip{D\partial_j\psi}{\partial_j\psi}=0$.
\end{proof}

\section{The identity on the original domain}

The manuscript proves the identity for $L^2$ eigenfunctions through the unitary group
$U_t$ generated by $D$. We use instead an exact algebraic identity and a mollifier that
smooths only across the orbits of $Y$. This requires $W$ to be supported away from the
separatrix and from the centers, which suffices for the even--even sector (the center is
cut out in Lemma~\ref{lem:ee-truncated-virial}).

\begin{lemma}[An exact identity]
\label{lem:virial-exact-identity}
\uses{def:eigen-problem}
\lean{NoEmbeddedEigenvalues.inner_virial_exact}
\leanok
Let $a=\omega+\psi$ and $Ha=\lambda\omega$ with $\lambda$ real. For every linear map $V$
defined on $\omega$ and $\psi$,
\[
  \ip{Va}{Ha}-\ip{VHa}{a}=\lambda\bigl(\ip{V\psi}{\omega}-\ip{V\omega}{\psi}\bigr).
\]
\end{lemma}

\begin{proof}
\leanok
Substitute $Ha=\lambda\omega$ and expand; the two terms $\lambda\ip{V\omega}{\omega}$ cancel.
\end{proof}

\begin{remark}
For $V=D$ and smooth functions, the left side is $\ip{[H,D]a}{a}$ and the right side is
$2\lambda\operatorname{Re}\ip{D\psi}{\omega}$, which is the proof of
Lemma~\ref{lem:virial-smooth}. The lemma replaces the manuscript's identity
$\ip{[H,U_t]a}{a}=\lambda\ip{[U_t,A]\omega}{\omega}$: neither $A=I+\Delta^{-1}$, nor its
extension to constants, nor the boundedness of $[D,A]$ is needed.
\end{remark}

\begin{lemma}[The smooth identity for $H^1$ functions]
\label{lem:virial-h1}
\uses{def:adapted-field, def:half-density-generator, def:sobolev-torus}
Let $W\in C^2$ be periodic and adapted with periodic coefficient $f\in C^1$. For every
$u\in H^1(\T^2)$, $\ip{Du}{v}=-\ip{u}{Dv}$ for $v\in H^1$, and
$2\operatorname{Re}\ip{Du}{Hu}=\operatorname{Re}\ip{fHu}{u}$.
\end{lemma}

\begin{proof}
\uses{lem:commutator-pointwise, lem:commutator-algebra, lem:sobolev-torus-basic}
For smooth $u$, $2\operatorname{Re}\ip{Du}{Hu}=\ip{HDu}{u}-\ip{DHu}{u}=\ip{[H,D]u}{u}
=\operatorname{Re}\ip{fHu}{u}$ by skew-symmetry and
Lemma~\ref{lem:commutator-pointwise}(3). Both sides are continuous in $H^1$, and Fourier
truncations converge in $H^1$.
\end{proof}

\begin{lemma}[A transversal mollifier]
\label{lem:transversal-mollifier}
\uses{def:action-angle, def:sobolev-torus}
Let $0<h_1'<h_1\le h_2<h_2'<1$, $\chi(h)$ smooth with support in $(h_1',h_2')$ and equal to
$1$ near $[h_1,h_2]$, and $\rho_\epsilon$ a standard mollifier. In the action-angle
coordinates of Definition~\ref{def:action-angle} let $J_\epsilon g=\rho_\epsilon*_h(\chi g)$,
convolution in $h$ at fixed $\theta$, extended by zero. Then for $\epsilon$ small:
\begin{enumerate}
  \item $J_\epsilon g\to\chi g$ in $L^2$ for $g\in L^2$, and in $H^1$ for $g\in H^1$;
  \item if $g,L_0g\in L^2$ then $J_\epsilon g\in H^1$ and $\epsilon\|\nabla J_\epsilon g\|\to0$;
  \item if $g,L_0g\in L^2$ then $\|[H,J_\epsilon]g\|\le C\epsilon\,(\|g\|+\|L_0g\|)$.
\end{enumerate}
\end{lemma}

\begin{proof}
\uses{lem:action-angle-chart}
In the chart $L_0=s(h)\partial_\theta$ with $s=1/S$, and $J_\epsilon$ commutes with
$\partial_h$, $\partial_\theta$ and with multiplication by functions of $h$ up to the
commutator below. (1) is the approximate identity property. For (2),
$\partial_\theta J_\epsilon g=J_\epsilon(\partial_\theta(\chi g))$ with
$\partial_\theta(\chi g)=\chi S\,L_0g\in L^2$, and
$\epsilon\partial_hJ_\epsilon g=(\rho')_\epsilon*_h(\chi g)\to0$ since $\int\rho'=0$. For (3),
\[
  [L_0,J_\epsilon]g(h,\theta)=\int\rho_\epsilon(h-h')\bigl(s(h)-s(h')\bigr)
    \partial_\theta(\chi g)(h',\theta)\,dh',
\]
and $|s(h)-s(h')|\le\epsilon\sup|s'|$ on the window. The area density $2S(h)$ is bounded
above and below on $[h_1',h_2']$, so the coordinate norms are equivalent to the $L^2$ norms.
\end{proof}

\begin{lemma}[Virial identity on the original domain]
\label{lem:virial-l2}
\uses{def:adapted-field, def:eigen-problem, def:half-density-generator, def:square-integral, def:stream-lifts}
\lean{NoEmbeddedEigenvalues.virial_l2}
\leanok
Let $0<h_1\le h_2<1$ and let $W$ be a periodic $C^3$ vector field, adapted with periodic
coefficient $f\in C^2$, with $W$ and $f$ supported in $\{2h_1\le\Psi\le2h_2\}$ (a compact
annulus in each positive cell, away from the separatrix and the centers). Then the
conclusions of Lemmas~\ref{lem:virial-smooth} and~\ref{lem:virial-smooth-quadratic} hold
for every eigenfunction $\omega\in X$ with $\lambda\ne0$ (no regularity beyond
Definition~\ref{def:eigen-problem}), with $\nabla\psi$ the Fourier-defined derivatives and
$\Delta\psi=\omega$ (manuscript \texttt{(ee:bounded-virial)}, restricted to such fields).
\end{lemma}

\begin{proof}
\uses{lem:virial-exact-identity, lem:virial-h1, lem:transversal-mollifier,
  lem:virial-smooth-quadratic, lem:elliptic-torus}
Here $\omega,L_0\omega\in L^2$ (since $L_0\omega=i\lambda\omega-L_0\psi$ and $\psi\in H^2$),
$\psi\in H^2$, $a=\omega+\psi$ and $Ha=\lambda\omega$. Take $J_\epsilon$ from
Lemma~\ref{lem:transversal-mollifier} with $\chi=1$ near the support of $W$, and apply
Lemma~\ref{lem:virial-exact-identity} with $V=J_\epsilon^\dagger DJ_\epsilon$
($\dagger$ the $L^2$ adjoint). The functions $J_\epsilon a$, $J_\epsilon\omega$, $J_\epsilon\psi$ are in
$H^1$. Right side: by skew-symmetry of $D$ on $H^1$,
$\ip{V\psi}{\omega}-\ip{V\omega}{\psi}=2\operatorname{Re}\ip{DJ_\epsilon\psi}{J_\epsilon\omega}
\to2\operatorname{Re}\ip{D\psi}{\omega}$, using $J_\epsilon\psi\to\chi\psi$ in $H^1$ and
$\chi D\psi=D\psi$. Left side: it equals $\ip{DJ_\epsilon a}{J_\epsilon Ha}-\ip{DJ_\epsilon Ha}{J_\epsilon a}$.
Write $J_\epsilon Ha=HJ_\epsilon a-[H,J_\epsilon]a$ and use skew-symmetry for the second term
($Ha=\lambda\omega$ has $L_0Ha\in L^2$, so $J_\epsilon Ha\in H^1$):
\[
  \text{left side}=2\operatorname{Re}\ip{DJ_\epsilon a}{HJ_\epsilon a}
    -2\operatorname{Re}\ip{DJ_\epsilon a}{[H,J_\epsilon]a}
  =\operatorname{Re}\ip{fHJ_\epsilon a}{J_\epsilon a}
    -2\operatorname{Re}\ip{DJ_\epsilon a}{[H,J_\epsilon]a}
\]
by Lemma~\ref{lem:virial-h1}. The first term tends to
$\operatorname{Re}\ip{fHa}{a}=\lambda\operatorname{Re}\ip{f\omega}{a}$. The second tends to
zero: $\|DJ_\epsilon a\|=o(1/\epsilon)$ and $\|[H,J_\epsilon]a\|=O(\epsilon)$ by
Lemma~\ref{lem:transversal-mollifier}. Divide by $\lambda$. The quadratic form follows by
the integrations by parts of Lemma~\ref{lem:virial-smooth-quadratic}, extended to
$\psi\in H^2$ by density.
\end{proof}

\begin{remark}[Why the mollifier smooths only across the orbits]
$H$ has no derivative across the orbits, so smoothing in $h$ almost commutes with it: the
commutator only sees the variation of the frequency $1/S(h)$ over the smoothing window,
which costs $O(\epsilon)$ times a $\theta$-derivative, and that derivative is controlled by
$H$ itself. With an ordinary mollifier in $x,y$ the pairing
$\operatorname{Re}\ip{DJ_\epsilon a}{[H,J_\epsilon]a}$ is only $O(1)$: its principal symbol does
not vanish on the directions conormal to the orbits, where an $L^2$ eigenfunction can be
rough. The flow has genuine twist ($S'\ne0$; numerically $S$ decreases from $\infty$ at the
separatrix to $2\pi$ at the center), so the commutator is not zero, only small.
\end{remark}

\begin{remark}[Digest note: the manuscript's argument]
The manuscript uses the flow $\Phi_t$ of $W$ and
$U_tu=\Jac(\Phi_t)^{1/2}u\circ\Phi_t$, a strongly continuous unitary group with generator
$D$, shows $HU_t=c_tU_tH-\tfrac i2(L_0\log c_t)U_t$ with
$c_t=\exp\int_0^tf\circ\Phi_\tau\,d\tau$, hence $[H,U_t]/t\to\tfrac12(fH+Hf)$ on the domain
of $H$, and passes to the limit in $\ip{[H,U_t]a}{a}=\lambda\ip{[U_t,A]\omega}{\omega}$.
The formulas check out (October 2026): $D\Phi_tY$ and $c_tY\circ\Phi_t$ both solve
$\dot v=DW(\Phi_t)v$ because $[Y,W]=fY$, and $L_0\log\Jac(\Phi_t)=L_0\log c_t$ because
$\Div W-f$ is constant on orbits. But the argument needs the Jacobian of $\Phi_t$
(differentiable dependence of the flow on its initial point, not in Mathlib), the group
properties of $U_t$, the core property of Lemma~\ref{lem:H-selfadjoint}, and, for fields
reaching the center, $C^{2,1}$ regularity there.
The proof above avoids all of these; the price is the action-angle chart
(Lemma~\ref{lem:action-angle-chart}) and the restriction to fields supported in an
annulus. Whether the odd--odd weighted analogue (Proposition~\ref{prop:dd-virial}) can be
treated in the same way is open.
\end{remark}
